% GENERATED FILE. Edit canonical lessons, then run npm run generate:books.
% canonical-source-hash: fnv1a64:e522c62aa56b19a1
% canonical-lessons: BN-C239-chumu-khaoya, BN-C239-rong-kora, BN-C239-bajano, BN-C239-phiriye-deoya, BN-C239-mone-rakha

\chapter{To kiss, To paint, To play music, To give back, To remember}
\label{ch:bn-to-kiss-to-paint-to-play-music-to-give-back-to-remember}
\hlchaptermodality{239}

\begin{chapteropening}
\textbf{By the end of this chapter:} \emph{I can say to kiss, to paint, to play music, to give back and to remember.}

Complete the last lesson of chapter 239: I can say to kiss, to paint, to play music, to give back and to remember.
\end{chapteropening}

\section[\texorpdfstring{\bn{চুমু} \bn{খাওয়া} (chumu khāoyā)}{চুমু খাওয়া (chumu khāoyā)}]{\bn{চুমু} \bn{খাওয়া} (chumu khāoyā) — to kiss}

\label{lesson:BN-C239-chumu-khaoya}



\begin{quote}
Before the new one: say the Bengali for to chase, then the Bengali for to believe.
\end{quote}

\subsection*{You'll want to know: \bn{চুমু} \bn{খাওয়া}}
\textbf{\bn{চুমু} \bn{খাওয়া}} — \emph{chumu khāoyā} — \textquotedblleft{}to kiss\textquotedblright{}.

\begin{grammarlens}[title={What you've built}]
One more word for this chapter.
\end{grammarlens}

\subsection*{Your turn}
\begin{itemize}
\raggedright
  \item \emph{Say it:} \emph{chumu khāoyā}
  \item \emph{Say it:} \emph{chumu khāoyā}, once more
  \item \emph{Say it:} say \emph{chumu khāoyā}
  \item \emph{Recall:} say the Bengali for to chase, then the Bengali for to believe, then say \emph{chumu khāoyā} again
\end{itemize}

\begin{tcolorbox}[breakable,colback=teal!4,colframe=teal!35!black,title={Before you move on}]
What does \bn{চুমু} \bn{খাওয়া} mean? (\textquotedblleft{}To kiss\textquotedblright{}.) Say it once more.
\end{tcolorbox}
\section[\texorpdfstring{\bn{রং} \bn{করা} (rông kôrā)}{রং করা (rông kôrā)}]{\bn{রং} \bn{করা} (rông kôrā) — to paint}

\label{lesson:BN-C239-rong-kora}



\begin{quote}
Before the new one: say the Bengali for to believe, then the Bengali for to kiss.
\end{quote}

\subsection*{You'll want to know: \bn{রং} \bn{করা}}
\textbf{\bn{রং} \bn{করা}} — \emph{rông kôrā} — \textquotedblleft{}to paint\textquotedblright{}.

\begin{grammarlens}[title={What you've built}]
One more word for this chapter.
\end{grammarlens}

\subsection*{Your turn}
\begin{itemize}
\raggedright
  \item \emph{Say it:} \emph{rông kôrā}
  \item \emph{Say it:} \emph{rông kôrā}, once more
  \item \emph{Say it:} say \emph{rông kôrā}
  \item \emph{Recall:} say the Bengali for to believe, then the Bengali for to kiss, then say \emph{rông kôrā} again
\end{itemize}

\begin{tcolorbox}[breakable,colback=teal!4,colframe=teal!35!black,title={Before you move on}]
What does \bn{রং} \bn{করা} mean? (\textquotedblleft{}To paint\textquotedblright{}.) Say it once more.
\end{tcolorbox}
\section[\texorpdfstring{\bn{বাজানো} (bājāno)}{বাজানো (bājāno)}]{\bn{বাজানো} (bājāno) — to play music, to play (an instrument)}

\label{lesson:BN-C239-bajano}



\begin{quote}
Before the new one: say the Bengali for to kiss, then the Bengali for to paint.
\end{quote}

\subsection*{You'll want to know: \bn{বাজানো}}
\textbf{\bn{বাজানো}} — \emph{bājāno} — \textquotedblleft{}to play music, to play (an instrument)\textquotedblright{}.

\begin{grammarlens}[title={What you've built}]
One more word for this chapter.
\end{grammarlens}

\subsection*{Your turn}
\begin{itemize}
\raggedright
  \item \emph{Say it:} \emph{bājāno}
  \item \emph{Say it:} \emph{bājāno}, once more
  \item \emph{Say it:} say \emph{bājāno}
  \item \emph{Recall:} say the Bengali for to kiss, then the Bengali for to paint, then say \emph{bājāno} again
\end{itemize}

\begin{tcolorbox}[breakable,colback=teal!4,colframe=teal!35!black,title={Before you move on}]
What does \bn{বাজানো} mean? (\textquotedblleft{}To play music, to play (an instrument)\textquotedblright{}.) Say it once more.
\end{tcolorbox}
\section[\texorpdfstring{\bn{ফিরিয়ে} \bn{দেওয়া} (phiriye deoyā)}{ফিরিয়ে দেওয়া (phiriye deoyā)}]{\bn{ফিরিয়ে} \bn{দেওয়া} (phiriye deoyā) — to give back}

\label{lesson:BN-C239-phiriye-deoya}



\begin{quote}
Before the new one: say the Bengali for to paint, then the Bengali for to play music.
\end{quote}

\subsection*{You'll want to know: \bn{ফিরিয়ে} \bn{দেওয়া}}
\textbf{\bn{ফিরিয়ে} \bn{দেওয়া}} — \emph{phiriye deoyā} — \textquotedblleft{}to give back\textquotedblright{}.

\begin{grammarlens}[title={What you've built}]
One more word for this chapter.
\end{grammarlens}

\subsection*{Your turn}
\begin{itemize}
\raggedright
  \item \emph{Say it:} \emph{phiriye deoyā}
  \item \emph{Say it:} \emph{phiriye deoyā}, once more
  \item \emph{Say it:} say \emph{phiriye deoyā}
  \item \emph{Recall:} say the Bengali for to paint, then the Bengali for to play music, then say \emph{phiriye deoyā} again
\end{itemize}

\begin{tcolorbox}[breakable,colback=teal!4,colframe=teal!35!black,title={Before you move on}]
What does \bn{ফিরিয়ে} \bn{দেওয়া} mean? (\textquotedblleft{}To give back\textquotedblright{}.) Say it once more.
\end{tcolorbox}
\section[\texorpdfstring{\bn{মনে} \bn{রাখা} (mône rākhā)}{মনে রাখা (mône rākhā)}]{\bn{মনে} \bn{রাখা} (mône rākhā) — to remember}

\label{lesson:BN-C239-mone-rakha}



\begin{quote}
Before the new one: say the Bengali for to play music, then the Bengali for to give back.
\end{quote}

\subsection*{You'll want to know: \bn{মনে} \bn{রাখা}}
\textbf{\bn{মনে} \bn{রাখা}} — \emph{mône rākhā} — \textquotedblleft{}to remember\textquotedblright{}.

\begin{grammarlens}[title={What you've built}]
One more word for this chapter.
\end{grammarlens}

\subsection*{Your turn}
\begin{itemize}
\raggedright
  \item \emph{Say it:} \emph{mône rākhā}
  \item \emph{Say it:} \emph{mône rākhā}, once more
  \item \emph{Say it:} say \emph{mône rākhā}
  \item \emph{Recall:} say the Bengali for to play music, then the Bengali for to give back, then say \emph{mône rākhā} again
\end{itemize}

\begin{tcolorbox}[breakable,colback=teal!4,colframe=teal!35!black,title={Before you move on}]
What does \bn{মনে} \bn{রাখা} mean? (\textquotedblleft{}To remember\textquotedblright{}.) Say it once more.
\end{tcolorbox}
